%% 6-8 2乗してマイナスになる数：虚数（第14回：組 V6-5。ページ制限なし。偶数ページ）
%% ほかの巻が指している形：2-8・3-10「2乗して −4 になる数は 6-8 で新しく考えます」→ x^2=-4 の解 ±2j（⑨）
%%   2-2「×(−1) で向きが反対」→「j をかけると90°回る」（②）。5-17「2乗すると −1 になる数 j（6-8）」、j^3=-j、j^4=1（④）
%%   3-10 の2つの特性方程式 s^2+3s+2=0、s^2+2s+5=0（解だけ。動きは 6-12）→ s=-1±2j（⑨）
\makeatletter
% 小見出し（太字の1行）  \bkviihead{見出し}
\newcommand\bkviihead[1]{\par\addvspace{6pt}\Needspace{6\baselineskip}\noindent{\sffamily\bfseries\color{\bkTT}#1}\par\nobreak\vspace{2pt}}
% 囲みの中の小見出し（それだけで1行）
\newcommand\bkviibox[1]{\par\addvspace{3pt}\noindent\textbf{#1}\par\noindent\ignorespaces}
% 大事な式のわく（巻の色のうすい地に、濃い色の線）  \bkviiframe[はば]{中身}
\newcommand\bkviiframe[2][100mm]{\par\addvspace{4pt}{\centering\tikz\node[draw=\bkTT, line width=0.9pt, fill=bkTheme!10, rounded corners=1.5mm, inner xsep=6mm, inner ysep=2.5mm, align=center]{#2};\par}\addvspace{4pt}\noindent\ignorespaces}
% 式を、まん中に1行で置く
\newcommand\bkviiline[1]{\par\addvspace{3pt}{\centering #1\par}\addvspace{3pt}\noindent\ignorespaces}
% 例題の途中式の行（折り返した行も、そろえた位置から始まる）
\newcommand\bkviiq{\par\vspace{2pt}\noindent\hangindent=4.3\zw\hangafter=1\hspace*{2.6\zw}}% （1）などで始まる行
\newcommand\bkviia{\par\vspace{2pt}\noindent\hangindent=2.6\zw\hangafter=0\relax}
\newcommand\bkviib{\par\vspace{2pt}\noindent\hangindent=4.3\zw\hangafter=0\relax}
% 平面の軸（原点中心）  \bkviiaxes{左}{右}{下}{上}（単位はその図の目もり）
\newcommand\bkviiaxes[4]{%
  \draw[bkGray, line width=0.6pt, -{Stealth[length=4pt]}] (#1,0) -- (#2,0);
  \draw[bkGray, line width=0.6pt, -{Stealth[length=4pt]}] (0,#3) -- (0,#4);}
% 点と矢印  \bkviiarrow{色}{(x,y)}
\newcommand\bkviiarrow[2]{\draw[#1, line width=1.5pt, -{Stealth[length=5pt]}, shorten >=0.4pt] (0,0) -- #2;}
\newcommand\bkviiit[2][0.24]{\makebox[#1\linewidth][l]{#2}}% 章末問題の小問を、幅をそろえて横に並べる（字の間はのびない）
\makeatother

\bkchapter{6}{8}{2乗してマイナスになる数：\\虚数}{\bklevel{4}}{2乗してマイナスになる数：虚数}
\bkmichishirube{2-2,2-8,3-8,3-10,4-8,5-17}{6-9,6-10,6-12}

\begin{goal}
\begin{itemize}[label={\color{\bkTT}●}, leftmargin=1.6\zw, itemsep=0pt, topsep=1pt]
\item \bklevel{4}\enspace 「$j$ をかけると90°回る」と考えて、$j^2=-1$ になるわけを説明できる。
\item \bklevel{4}\enspace 複素数のたし算・ひき算・かけ算・わり算ができる。
\item \bklevel{4}\enspace 判別式がマイナスの2次方程式を解いて、虚数の解を求められる。
\end{itemize}
\end{goal}

\begin{hatena}
\noindent\begin{minipage}[c]{0.6\linewidth}
\raggedright
2-2で、$\times(-1)$ は「向きを反対にする」かけ算だと学びました。数直線の上で3に $-1$ をかけると、$-3$ になります。0を中心に、ぐるっと半周（180°）回ったのと同じです。もう1回かけると、3にもどります。

では、「2回くり返すと $\times(-1)$ になる」かけ算はあるでしょうか。それは、2乗すると $-1$ になる数をかけることです。そんな数は、どこにあるのでしょう。
\end{minipage}\hfill
\begin{minipage}[c]{0.37\linewidth}
\centering
\begin{tikzpicture}[x=4.6mm, y=4.6mm]
  \draw[bkInk, line width=0.8pt, -{Stealth[length=4pt]}] (-4.3,0) -- (4.5,0);
  \foreach \k in {-3,...,3} {\draw[bkInk, line width=0.6pt] (\k,-0.15) -- (\k,0.15);}
  \node[font=\sffamily\footnotesize, anchor=north] at (-3,-0.35) {$-3$};
  \node[font=\sffamily\footnotesize, anchor=north] at (0,-0.35) {0};
  \node[font=\sffamily\footnotesize, anchor=north] at (3,-0.35) {3};
  \draw[bkOrange, line width=1.1pt, -{Stealth[length=4.5pt]}] (5:3) arc (5:86:3);
  \draw[bkOrange, line width=1.1pt, -{Stealth[length=4.5pt]}] (94:3) arc (94:175:3);
  \fill[bkOrange] (3,0) circle (0.2);
  \fill[bkBlue] (-3,0) circle (0.2);
  \node[font=\sffamily\bfseries\small, text=bkOrange!80!black] at (45:3.9) {？};
  \node[font=\sffamily\bfseries\small, text=bkOrange!80!black] at (135:3.9) {？};
  \node[font=\sffamily\footnotesize, anchor=south] at (0,3.35) {2回で $\times(-1)$};
\end{tikzpicture}
\end{minipage}
\end{hatena}

\noindent 3-10では、$x^2=-4$ のような方程式について「解はありません（このような数は 6-8 で新しく考えます）」と書きました。この章で、その約束を果たします。新しく出てくる数は、なぞの数ではありません。「90°回す」という、はっきりした意味をもつ数です。電気の回路やゆれの計算（6-10〜6-13）では、この数を毎日のように使います。

%% ---------------- やってみよう・思い出そう ----------------
\begin{yattemiyou}[方眼紙の上で、90°ずつ回そう]
\textbf{用意するもの}\par
方眼紙、定規、分度器
\begin{enumerate}[label={\bkcircled[bkGreen]{\arabic*}}, leftmargin=2.4\zw, itemsep=1pt, topsep=3pt]
\item 方眼紙のまん中に原点Oをとり、横の線とたての線を引こう。Oから右へ3目もりの点をAとし、OからAへ矢印をかこう。
\item 矢印OAを、Oを中心に左回りに90°回そう。矢印の先の点をBとしよう。
\item 矢印OBを、さらに左回りに90°回そう。矢印の先の点をCとしよう。
\item 点P $(3,\ 1)$ でも、②③と同じように2回回して、点Q、点Rをかこう。
\end{enumerate}
\begin{bkcenter}
\begin{bktable}{colspec={Q[c,m,wd=30mm] Q[c,m,wd=30mm] Q[c,m,wd=30mm]}, rowsep=1.5pt}
はじめの点 & 90°回す & もう90°回す \\
A $(3,\ 0)$ & B $(0,\ 3)$ & C $(-3,\ 0)$ \\
P $(3,\ 1)$ & Q $(-1,\ 3)$ & R $(-3,\ -1)$ \\
\end{bktable}
\bktabcaption{ミナさんがかいた結果}
\end{bkcenter}
\end{yattemiyou}

\begin{bkfigure}
\begin{tikzpicture}[x=7mm, y=7mm]
  \begin{scope}
    \draw[bkLine, line width=0.4pt, step=1] (-3.6,-1.6) grid (3.6,3.6);
    \bkviiaxes{-3.8}{3.9}{-1.8}{3.9}
    \draw[bkGray, line width=0.8pt, dash pattern=on 2pt off 1.5pt, -{Stealth[length=4pt]}] (8:1.6) arc (8:82:1.6);
    \draw[bkGray, line width=0.8pt, dash pattern=on 2pt off 1.5pt, -{Stealth[length=4pt]}] (98:1.6) arc (98:172:1.6);
    \bkviiarrow{bkBlue}{(3,0)}
    \bkviiarrow{bkOrange}{(0,3)}
    \bkviiarrow{bkRed}{(-3,0)}
    \node[font=\sffamily\footnotesize, anchor=north, inner sep=1pt] at (3,-0.3) {A $(3,\ 0)$};
    \node[font=\sffamily\footnotesize, anchor=west, inner sep=1pt, fill=white] at (0.3,3.1) {B $(0,\ 3)$};
    \node[font=\sffamily\footnotesize, anchor=north, inner sep=1pt] at (-3,-0.3) {C $(-3,\ 0)$};
    \node[font=\sffamily\footnotesize, anchor=north east, inner sep=1pt, fill=white] at (-0.15,-0.15) {O};
  \end{scope}
  \begin{scope}[xshift=62mm]
    \draw[bkLine, line width=0.4pt, step=1] (-3.6,-1.6) grid (3.6,3.6);
    \bkviiaxes{-3.8}{3.9}{-1.8}{3.9}
    \draw[bkGray, line width=0.8pt, dash pattern=on 2pt off 1.5pt, -{Stealth[length=4pt]}] (26:1.7) arc (26:101:1.7);
    \draw[bkGray, line width=0.8pt, dash pattern=on 2pt off 1.5pt, -{Stealth[length=4pt]}] (116:1.7) arc (116:191:1.7);
    \bkviiarrow{bkBlue}{(3,1)}
    \bkviiarrow{bkOrange}{(-1,3)}
    \bkviiarrow{bkRed}{(-3,-1)}
    \node[font=\sffamily\footnotesize, anchor=west, inner sep=1pt, fill=white] at (3.15,1.35) {P};
    \node[font=\sffamily\footnotesize, anchor=south, inner sep=1pt, fill=white] at (-1,3.15) {Q};
    \node[font=\sffamily\footnotesize, anchor=north, inner sep=1pt] at (-3,-1.2) {R};
    \node[font=\sffamily\footnotesize, anchor=north east, inner sep=1pt, fill=white] at (-0.15,-0.15) {O};
  \end{scope}
\end{tikzpicture}
\bkfigcaption{左：Aを90°ずつ2回回すと、Cは $-3$。右：Pを2回回したRは、Pと原点について反対の位置。}
\end{bkfigure}

\noindent Cは $-3$ で、Aの3に $-1$ をかけた数です。R $(-3,\ -1)$ も、P $(3,\ 1)$ の横と高さの両方に $-1$ をかけたものです。90°回すことを2回くり返すと、180°回すことになり、$\times(-1)$ と同じはたらきになりました。4-8の「90°回すと $(x,\ y)\to(-y,\ x)$」を2回使っても、
\bkviiline{$(x,\ y)\ \to\ (-y,\ x)\ \to\ (-x,\ -y)$}
で、同じことがわかります。2回目は、1回目の答え $(-y,\ x)$ の横 $-y$ と高さ $x$ を、もう一度 $(x,\ y)\to(-y,\ x)$ に入れました。

\begin{omoidasou}
\bkrecall{2-2}{$\times(-1)$ で向きが反対になる。2回かけると、もとの向きにもどる。}
\bkrecall{2-8}{正の数も負の数も、2乗すると正。$\sqrt a\times\sqrt b=\sqrt{ab}$（$a$、$b$ は正の数）。}
\bkrecall{2-12}{文字の部分が同じ項（同類項）どうしは、まとめられる：$3a+a=4a$}
\bkrecall{3-8}{$(a+b)(c+d)=ac+ad+bc+bd$、$(A+B)(A-B)=A^2-B^2$}
\bkrecall{3-10}{$ax^2+bx+c=0$ の解は $x=\dfrac{-b\pm\sqrt{b^2-4ac}}{2a}$。判別式 $D=b^2-4ac$。\\分母と分子に同じ数をかけて、分母の $\sqrt{\phantom{2}}$ をなくす（有理化）。}
\bkrecall{4-8}{原点のまわりに左回りに90°回すと $(x,\ y)\to(-y,\ x)$。}
\bkrecall{4-11}{ベクトルのたし算は、横どうし・たてどうしをたす（成分ごと）。}
\end{omoidasou}

%% ---------------- 数学の道具 ----------------
\bksection{doguu}{数学の道具}

\bksubsection{1}{2乗してマイナスになる数は、数直線の上にない}

\leavevmode\label{bk:6-8:nai}はじめに、なぜ「2乗して $-4$ になる数はない」と言ってきたのかを、もう一度確かめます。

\begin{bkcenter}
\begin{bktable}{colspec={Q[c,m,wd=14mm] *{7}{Q[c,m,wd=10mm]}}, rowsep=1pt}
$x$ & $-3$ & $-2$ & $-1$ & 0 & 1 & 2 & 3 \\
$x^2$ & 9 & 4 & 1 & 0 & 1 & 4 & 9 \\
\end{bktable}
\end{bkcenter}
\noindent\begin{minipage}[c]{0.46\linewidth}
\raggedright
表の下の段は、どれも0以上です。正の数どうしをかけても、負の数どうしをかけても（2-2）、答えは正になるからです。グラフ $y=x^2$（3-9）も、横軸より下にはきません。
\end{minipage}\hfill
\begin{minipage}[c]{0.5\linewidth}
\centering
\begin{tikzpicture}
\begin{axis}[bkplot, width=62mm, height=46mm, xmin=-3.3, xmax=3.3, ymin=-5.6, ymax=10, xtick={-3,-2,-1,1,2,3}, ytick={4,8}, xlabel={$x$}, ylabel={$y$},
    yticklabel style={fill=white, inner sep=1pt, outer sep=2.3pt}]
  \addplot[bkcurve, domain=-3.1:3.1, samples=60] {x^2};
  \addplot[bkRed, line width=1.1pt, dash pattern=on 4pt off 2pt, domain=-3.1:3.1] {-4};
  \node[font=\sffamily\footnotesize, text=bkBlue, anchor=west, fill=white, inner sep=1pt] at (axis cs:1.75,8.6) {$y=x^2$};
  \node[font=\sffamily\footnotesize, text=bkRed!80!black, anchor=south, fill=white, inner sep=1pt] at (axis cs:0.4,-3.7) {$y=-4$ の線とは交わらない};
\end{axis}
\end{tikzpicture}
\end{minipage}

\par\vspace{4pt}
これまでの数（数直線の上にならぶ数）を、まとめて実数といいました（4-10）。実数を2乗すると、かならず0以上になります。だから、2乗して $-4$ になる数は、数直線のどこをさがしても見つかりません。

\bkviihead{数の世界は、これまでも広がってきた}
答えのない計算に出会ったとき、わたしたちはこれまでも、新しい数を考えてきました。
\begin{bkcenter}
\begin{bktable}{colspec={Q[l,m,wd=46mm] Q[l,m,wd=50mm] Q[c,m,wd=22mm]}, rowsep=1.5pt}
答えのなかった計算 & 新しく考えた数 & 章 \\
$3-5$ & 負の数 $-2$ & 2-2 \\
$2\div3$ & 分数 $\dfrac{2}{3}$ & 1-16・2-4 \\
2乗して2になる数 & $\sqrt2$（分数で書けない数） & 2-8・2-9 \\
2乗して $-4$ になる数 & ？（この章） & 6-8 \\
\end{bktable}
\bktabcaption{答えのない計算に、新しい数で答えてきた。}
\end{bkcenter}
負の数も、はじめは「0より小さい数なんて、ありえない」と考えられていました。それでも、温度や借金や向きを表すのにとても便利なので、使われるようになりました。2乗してマイナスになる数も、同じです。数直線の上にはないけれど、新しく考えると、とても便利なのです。どこに置けばよいかを、次の②で考えます。

\bksubsection{2}{90°回すかけ算：虚数単位 $j$}

\leavevmode\label{bk:6-8:kaiten}2-2では、$\times(-1)$ を「向きを反対にする」かけ算だと考えました。向きを反対にすることは、0を中心に180°回すことと同じです。

\begin{bkfigure}
\begin{tikzpicture}[x=15mm, y=15mm]
  \begin{scope}
    \draw[bkInk, line width=0.8pt, -{Stealth[length=4pt]}] (-1.7,0) -- (1.8,0);
    \foreach \k in {-1,0,1} {\draw[bkInk, line width=0.6pt] (\k,-0.06) -- (\k,0.06);}
    \node[font=\sffamily\footnotesize, anchor=north] at (-1,-0.1) {$-1$};
    \node[font=\sffamily\footnotesize, anchor=north] at (0,-0.1) {0};
    \node[font=\sffamily\footnotesize, anchor=north] at (1,-0.1) {1};
    \draw[bkOrange, line width=1.2pt, -{Stealth[length=5pt]}] (6:1) arc (6:174:1);
    \fill[bkBlue] (1,0) circle (0.055);
    \fill[bkRed] (-1,0) circle (0.055);
    \node[font=\sffamily\footnotesize, anchor=south, text=bkOrange!80!black] at (0,1.08) {$\times(-1)$：180°回す};
    \node[font=\sffamily\small, anchor=north] at (0,-0.55) {数直線の上だけで考える};
  \end{scope}
  \begin{scope}[xshift=62mm]
    \draw[bkGray, line width=0.6pt, -{Stealth[length=4pt]}] (-1.7,0) -- (1.8,0);
    \draw[bkGray, line width=0.6pt, -{Stealth[length=4pt]}] (0,-0.45) -- (0,1.55);
    \foreach \k in {-1,1} {\draw[bkInk, line width=0.6pt] (\k,-0.06) -- (\k,0.06);}
    \node[font=\sffamily\footnotesize, anchor=north] at (-1,-0.1) {$-1$};
    \node[font=\sffamily\footnotesize, anchor=north] at (1,-0.1) {1};
    \draw[bkPurple, line width=1.2pt, -{Stealth[length=5pt]}] (6:1) arc (6:84:1);
    \draw[bkPurple, line width=1.2pt, -{Stealth[length=5pt]}] (96:1) arc (96:174:1);
    \fill[bkBlue] (1,0) circle (0.055);
    \fill[bkPurple] (0,1) circle (0.055);
    \fill[bkRed] (-1,0) circle (0.055);
    \node[font=\sffamily\footnotesize, anchor=west, fill=white, inner sep=1pt] at (0.1,1.12) {$j$（上へ1）};
    \node[font=\sffamily\footnotesize, text=bkPurple] at (45:1.28) {$\times j$};
    \node[font=\sffamily\footnotesize, text=bkPurple] at (135:1.28) {$\times j$};
    \node[font=\sffamily\small, anchor=north] at (0,-0.55) {平面に広げて考える};
  \end{scope}
\end{tikzpicture}
\bkfigcaption{左：$\times(-1)$ は180°回す。右：$\times j$ は90°回す。2回で180°回って、1が $-1$ になる。}
\end{bkfigure}

\bkviihead{見方1：図で考える}
「2回くり返すと $\times(-1)$ になるかけ算」は、180°の半分、つまり90°回すかけ算です。ところが、数直線の上だけでは、90°回すことはできません。90°回すと、数直線の外に出てしまうからです。

そこで、数直線を平面に広げます。これまでの数（実数）は、横の線の上にならべます。そして、「0を中心に左回りに90°回す」かけ算を、$\times j$（ジェイをかける）と書くことにします。
\begin{itemize}[leftmargin=1.2\zw, itemsep=2pt, topsep=2pt, label={\color{\bkTT}●}]
\item 1に $j$ をかけると、1を90°回した所、つまり0から真上へ1の所に行きます。$1\times j=j$ なので、ここが数 $j$ そのものの場所です。
\item その $j$ に、もう一度 $j$ をかけると、さらに90°回って、$-1$ の所に行きます。
\end{itemize}
2回の $\times j$ で $\times(-1)$ になったので、
\bkviiframe[70mm]{\large $j\times j=-1$\qquad つまり\qquad $j^2=-1$}

\bkviihead{見方2：座標で計算する}
やってみようと同じく、4-8の「90°回すと $(x,\ y)\to(-y,\ x)$」で計算します。1は横の線の上の点 $(1,\ 0)$ です。
\bkviiline{$(1,\ 0)\ \xrightarrow{\ \times j\ }\ (0,\ 1)\ \xrightarrow{\ \times j\ }\ (-1,\ 0)$}
1回目で $(0,\ 1)$（真上へ1。これが $j$）、2回目で $(-1,\ 0)$（数直線の $-1$）。図で考えたことと、ぴったり合います。

\bkviihead{右回りでもよい：$-j$}
右回りに90°回しても、2回くり返すと180°回って $-1$ になります。右回りに90°回した所、つまり真下へ1の所の数は、$j$ の反対の $-j$ です（2-2の「反対の数」）。式で確かめると、
\begin{bkhenkei}
(-j)^2=(-1)^2\times j^2 & $-j=(-1)\times j$ と見て、2乗を分ける \\
=1\times(-1) & $(-1)^2=1$、$j^2=-1$ \\
=-1 & \\
\end{bkhenkei}
2乗して1になる数が1と $-1$ の2つあるように、2乗して $-1$ になる数も、$j$ と $-j$ の2つです。

\bkviihead{\bkhatenaref の答え}
2乗すると $-1$ になる数は、数直線の上にはありません。数直線を平面に広げたときの、\textbf{0から真上へ1の所}（$j$）と、\textbf{真下へ1の所}（$-j$）にあります。「2回くり返すと $\times(-1)$ になるかけ算」は、「90°回す」かけ算だったのです。

この $j$ を、\textbf{虚数単位}（きょすうたんい）といいます。\bkidx{きょすうたんい}{虚数単位}「虚」は「目に見える形がない」という意味で、昔の人が「想像の中だけの数」と考えたことからついた名前です。でも、この章で見るように、$j$ には「90°回す」というはっきりした意味があります。

\begin{seikaku}
数学では、$j$ を「2乗すると $-1$ になる新しい数」として決め（定義し）、そのうえで、たし算・かけ算の入れかえや、かっこのはずし方（分配法則）などの、これまでの計算のきまりが、そのまま使えるように数の世界を広げます。そう決めても計算にむじゅんが起きないことは、1830年代にハミルトンが、複素数（③）を「2つの実数の組」として決めなおして確かめました。「90°回す」という図は、そう決めた $j$ がどうふるまうかを目で見る、1つの見方です。
\end{seikaku}

\bkviihead{$i$ ではなく $j$ と書くわけ}
数学の本や高校の教科書では、虚数単位を $i$（imaginary unit の頭文字）と書きます。ところが電気では、$i$ を電流の文字に使います（電流の強さを表すフランス語 intensité の頭文字）。そこで電気や制御の本では、虚数単位を $j$ と書きます。この本は、電気や機械で使うための数学なので、$j$ を使います。プログラムの言語 Python も、$3+4j$ のように $j$ を使います（⑩）。$i$ と $j$ は、書き方がちがうだけで、同じ数です。

\bksubsection{3}{複素数 $a+bj$：実部と虚部}

\leavevmode\label{bk:6-8:fukuso}$j$ に実数をかけたり、実数をたしたりした数を考えます。

$3j$ は、1を90°回した $j$ を3倍した数なので、0から真上へ3の所です。$-2j$ は真下へ2の所です。そして、$3+4j$ は「右へ3進んで、上へ4進んだ所」、つまり平面の点 $(3,\ 4)$ です。

\begin{bkfigure}
\begin{tikzpicture}[x=6.2mm, y=6.2mm]
  \begin{scope}
    \draw[bkLine, line width=0.4pt, step=1] (-3.6,-2.6) grid (4.6,4.6);
    \bkviiaxes{-3.8}{4.9}{-2.8}{4.9}
    \node[font=\sffamily\footnotesize, anchor=north west, inner sep=1pt] at (4.3,-0.12) {実数};
    \foreach \x/\y/\c in {3/4/bkBlue, -2/1/bkOrange, 0/2/bkPurple, 2.5/0/bkGreen!70!black, -1/-2/bkRed} {\fill[\c] (\x,\y) circle (0.16);}
    \node[font=\sffamily\footnotesize, anchor=west, fill=white, inner sep=1pt] at (3.2,4) {$3+4j$};
    \node[font=\sffamily\footnotesize, anchor=south, fill=white, inner sep=1pt] at (-2,1.25) {$-2+j$};
    \node[font=\sffamily\footnotesize, anchor=west, fill=white, inner sep=1pt] at (0.25,2) {$2j$};
    \node[font=\sffamily\footnotesize, anchor=south, fill=white, inner sep=1pt] at (2.5,0.25) {2.5};
    \node[font=\sffamily\footnotesize, anchor=east, fill=white, inner sep=1pt] at (-1.25,-2) {$-1-2j$};
    \draw[bkBlue, line width=0.6pt, dash pattern=on 2pt off 1.5pt] (3,0) -- (3,4) (0,4) -- (3,4);
  \end{scope}
  % 数の仲間（入れ子の箱）
  \begin{scope}[x=1mm, y=1mm, xshift=44mm, yshift=-17mm]
    \draw[bkPurple, line width=0.9pt, rounded corners=2mm, fill=bkPurple!5] (0,0) rectangle (60,56);
    \draw[bkBlue, line width=0.9pt, rounded corners=2mm, fill=bkBlue!5] (2.5,2.5) rectangle (57.5,45);
    \draw[bkGreen!70!black, line width=0.9pt, rounded corners=2mm, fill=bkGreen!6] (5,5) rectangle (55,34);
    \draw[bkOrange, line width=0.9pt, rounded corners=2mm, fill=bkOrange!7] (7.5,7.5) rectangle (52.5,23);
    \draw[bkRed, line width=0.9pt, rounded corners=2mm, fill=bkRed!5] (10,10) rectangle (50,17);
    \node[font=\sffamily\bfseries\footnotesize, anchor=west, text=bkPurple] at (3,50.5) {複素数};
    \node[font=\sffamily\footnotesize, anchor=east] at (57,50.5) {$2j$、$3+4j$};
    \node[font=\sffamily\bfseries\footnotesize, anchor=west, text=bkBlue!80!black] at (5.5,39.5) {実数};
    \node[font=\sffamily\footnotesize, anchor=east] at (54.5,39.5) {$\sqrt2$、$\pi$};
    \node[font=\sffamily\bfseries\footnotesize, anchor=west, text=bkGreen!50!black] at (8,28.5) {有理数（分数）};
    \node[font=\sffamily\footnotesize, anchor=east] at (52,28.5) {$\frac{1}{2}$、0.3};
    \node[font=\sffamily\bfseries\footnotesize, anchor=west, text=bkOrange!80!black] at (10.5,20) {整数};
    \node[font=\sffamily\footnotesize, anchor=east] at (49.5,20) {0、$-3$};
    \node[font=\sffamily\bfseries\footnotesize, anchor=west, text=bkRed!80!black] at (12.5,13.5) {自然数};
    \node[font=\sffamily\footnotesize, anchor=east] at (47.5,13.5) {1、2、3};
  \end{scope}
\end{tikzpicture}
\bkfigcaption{左：複素数は平面の点。横の線の上が実数。右：数の仲間。外側の箱ほど広い。}\label{bk:6-8:fig-nakama}
\end{bkfigure}

\bkviiframe[110mm]{$a$、$b$ を実数として、$a+bj$ の形の数を\textbf{複素数}（ふくそすう）という。\\[2pt] $a$ を\textbf{実部}（じつぶ）、$b$ を\textbf{虚部}（きょぶ）という。}
\bkidx{ふくそすう}{複素数}\bkidx{じつぶ}{実部}\bkidx{きょぶ}{虚部}「複」は「いくつかが組になった」という意味です。複素数は、実部と虚部の2つの実数が組になった数です。複素数を1つの文字で表すときは、$z$（ゼット）をよく使います。$z$ の実部を $\mathrm{Re}\,z$、虚部を $\mathrm{Im}\,z$ と書きます（real part、imaginary part の頭文字）。たとえば $z=3+4j$ なら、
\bkviiline{$\mathrm{Re}\,z=3$\qquad $\mathrm{Im}\,z=4$}
虚部は $4j$ ではなく、$j$ にかかっている実数の4です。

\begin{bkcenter}
\begin{bktable}{colspec={Q[l,m,wd=24mm] Q[l,m,wd=52mm] Q[l,m,wd=40mm]}, rowsep=1.5pt}
名前 & どんな数か & 例 \\
実数 & 虚部が0（$b=0$） & $2.5$、$-1$、$\sqrt2$ \\
虚数 & 虚部が0でない（$b\neq0$） & $3+4j$、$-1-2j$、$2j$ \\
純虚数 & 実部が0の虚数（$a=0$） & $2j$、$-j$ \\
\end{bktable}
\bktabcaption{実数も、虚部が0の複素数と考える。}
\end{bkcenter}
実数 $2.5$ は $2.5+0j$ と考えられるので、実数も複素数の仲間です。図\ref{bk:6-8:fig-nakama}の右のように、数の仲間は、これまでも外側へ広がってきました。複素数は、そのいちばん外側です。

2つの複素数が等しいのは、平面の同じ点になるとき、つまり\textbf{実部どうしも、虚部どうしも等しいとき}です。たとえば $a+bj=3-2j$ なら、$a=3$、$b=-2$ です。

\bksubsection{4}{$j$ を何回もかける：4回でもとにもどる}

\leavevmode\label{bk:6-8:beki}$j$ を何回もかけた数（$j$ の累乗。2-7）を計算します。$j^2=-1$ を使います。

\noindent\begin{minipage}[c]{0.56\linewidth}
\begin{bktable}{colspec={Q[c,m,wd=8mm] Q[l,m,wd=42mm] Q[c,m,wd=10mm]}, rowsep=1pt, cells={font=\small}}
 & 計算 & 答え \\
$j^1$ & & $j$ \\
$j^2$ & 決めた数 & $-1$ \\
$j^3$ & $j^2\times j=(-1)\times j$ & $-j$ \\
$j^4$ & $j^3\times j=(-j)\times j=-j^2$ & 1 \\
$j^5$ & $j^4\times j=1\times j$ & $j$ \\
$j^6$ & $j^4\times j^2=1\times(-1)$ & $-1$ \\
\end{bktable}
\end{minipage}\hfill
\begin{minipage}[c]{0.4\linewidth}
\centering
\begin{tikzpicture}[x=15mm, y=15mm]
  \draw[bkGray, line width=0.6pt, -{Stealth[length=4pt]}] (-1.55,0) -- (1.6,0);
  \draw[bkGray, line width=0.6pt, -{Stealth[length=4pt]}] (0,-1.5) -- (0,1.55);
  \draw[bkInk!40, line width=0.7pt] (0,0) circle (1);
  \foreach \a in {0,90,180,270} {\draw[bkPurple, line width=1pt, -{Stealth[length=4pt]}] (\a+9:1) arc (\a+9:\a+81:1);}
  \foreach \a in {0,90,180,270} {\fill[bkInk] (\a:1) circle (0.05);}
  \node[font=\sffamily\footnotesize, anchor=south west, inner sep=1pt] at (1.04,0.05) {$1=j^4$};
  \node[font=\sffamily\footnotesize, anchor=south west, inner sep=1pt] at (0.06,1.04) {$j=j^5$};
  \node[font=\sffamily\footnotesize, anchor=south east, inner sep=1pt] at (-1.04,0.05) {$-1=j^2$};
  \node[font=\sffamily\footnotesize, anchor=north west, inner sep=1pt] at (0.06,-1.04) {$-j=j^3$};
\end{tikzpicture}
\end{minipage}

\par\vspace{4pt}
$j$ をかけるたびに90°ずつ回るので、4回で360°回って、もとの1にもどります。だから、$j$ の累乗は $j$、$-1$、$-j$、1 の4つをくり返します。$j^n$ を求めるには、$n$ を4でわったあまりを見ればよいのです。1-15の「100日後の曜日」や、4-8の「何周と何度」と同じ考え方です。たとえば $j^{10}$ は、$10=4\times2+2$ なので、4回ずつ2組かけると1にもどり、あまりの2回分だけが残って $j^{10}=j^2=-1$ です。

5-17のオイラーの公式への入口で使った $j^2=-1$、$j^3=-j$、$j^4=1$ は、この表の3行です。

\bksubsection{5}{たし算・ひき算：実部どうし、虚部どうし}

\leavevmode\label{bk:6-8:tasu}\noindent\begin{minipage}[c]{0.58\linewidth}
\raggedright
$j$ を、$x$ のような文字と同じように考えます。2-12で、$3+4x$ と $1-2x$ をたすときは、数どうし、$x$ のつく項どうし（同類項）をまとめました。複素数も同じです。
\bkviiline{$(3+4j)+(1-2j)$}\bkviiline{$=(3+1)+(4-2)j=4+2j$}
平面で見ると、「右へ3・上へ4」の矢印のあとに、「右へ1・下へ2」の矢印をつないだ所です。4-11のベクトルのたし算（成分ごとにたす）と同じです。
\end{minipage}\hfill
\begin{minipage}[c]{0.38\linewidth}
\centering
\begin{tikzpicture}[x=6mm, y=6mm]
  \draw[bkLine, line width=0.4pt, step=1] (-0.6,-0.6) grid (4.6,4.6);
  \bkviiaxes{-0.8}{4.9}{-0.8}{4.9}
  \bkviiarrow{bkBlue}{(3,4)}
  \draw[bkOrange, line width=1.5pt, -{Stealth[length=5pt]}] (3,4) -- (4,2);
  \draw[bkPurple, line width=1.5pt, -{Stealth[length=5pt]}] (0,0) -- (4,2);
  \node[font=\sffamily\footnotesize, anchor=east, fill=white, inner sep=1pt] at (2.7,4.1) {$3+4j$};
  \node[font=\sffamily\footnotesize, anchor=west, fill=white, inner sep=1pt] at (3.75,3.3) {$1-2j$};
  \node[font=\sffamily\footnotesize, anchor=north west, fill=white, inner sep=1pt] at (2.1,1.0) {$4+2j$};
\end{tikzpicture}
\end{minipage}

\par\vspace{4pt}
ひき算も同じで、$(3+4j)-(1-2j)=(3-1)+(4-(-2))j=2+6j$ です。かっこの前の $-$ は、かっこの中の両方の項につくことに注意しましょう。実数をかけるときは、分配法則で両方の項にかけます。$2(3+4j)=6+8j$。

\bksubsection{6}{かけ算：展開して、$j^2=-1$}

\leavevmode\label{bk:6-8:kakeru}かけ算も、$j$ を文字のように考えて、3-8の展開をします。そのあとで $j^2$ を $-1$ におきかえます。$(3+4j)(1-2j)$ を、1行ずつ計算します。まず、3-8の「全部の組み合わせをかけてたす」で、
\bkviiline{$(3+4j)(1-2j)=3\times1+3\times(-2j)+4j\times1+4j\times(-2j)$}
となります。ここから先を、1行に1つずつ進めます。
\begin{bkhenkei}
=3-6j+4j-8j^2 & それぞれのかけ算をする \\
=3-6j+4j-8\times(-1) & $j^2=-1$ に入れかえる \\
=3+8-2j & 実部どうし、虚部どうしを集める \\
=11-2j & \\
\end{bkhenkei}
$j^2$ が出てきたら、すぐに $-1$ にすることが大切です。このかけ算は、文字のままでも同じように計算できます。
\bkviiline{$(a+bj)(c+dj)=ac+adj+bcj+bdj^2=(ac-bd)+(ad+bc)j$}
この式は覚えなくてかまいません。展開して $j^2=-1$ を使えば、いつでも作れます。

\bkviihead{$j$ をかけると、ほんとうに90°回る？}
\noindent\begin{minipage}[c]{0.56\linewidth}
\raggedright
②では、「90°回すかけ算」を $\times j$ と決めました。では、$3+4j$ に $j$ をかけて、展開で計算してみます。
\bkviiline{$j(3+4j)=3j+4j^2=3j-4=-4+3j$}
点 $(3,\ 4)$ が、点 $(-4,\ 3)$ に移りました。4-8の「90°回すと $(x,\ y)\to(-y,\ x)$」で $x=3$、$y=4$ とすると $(-4,\ 3)$。ぴったり合います。展開のきまりで計算しても、$j$ をかけることは、たしかに90°回すことなのです。
\end{minipage}\hfill
\begin{minipage}[c]{0.4\linewidth}
\centering
\begin{tikzpicture}[x=5.6mm, y=5.6mm]
  \draw[bkLine, line width=0.4pt, step=1] (-4.6,-0.6) grid (3.6,4.6);
  \bkviiaxes{-4.8}{3.9}{-0.8}{4.9}
  \bkviiarrow{bkBlue}{(3,4)}
  \bkviiarrow{bkPurple}{(-4,3)}
  \draw[bkInk, line width=0.5pt] ($(0,0)!0.12!(3,4)$) -- ++($(0,0)!0.12!(-4,3)$) -- ($(0,0)!0.12!(-4,3)$);
  \draw[bkGray, line width=0.8pt, dash pattern=on 2pt off 1.5pt, -{Stealth[length=4pt]}] (57:2.6) arc (57:139:2.6);
  \node[font=\sffamily\footnotesize, anchor=west, fill=white, inner sep=1pt] at (3.1,4.1) {$3+4j$};
  \node[font=\sffamily\footnotesize, anchor=south, fill=white, inner sep=1pt] at (-3.2,3.25) {$-4+3j$};
  \node[font=\sffamily\footnotesize, text=bkGray, fill=white, inner sep=1pt] at (98:3.25) {90°};
\end{tikzpicture}
\end{minipage}

\par\vspace{4pt}
②では図から $j^2=-1$ を決め、ここでは $j^2=-1$ から計算して「90°回る」を確かめました。2つの見方は、行ったり来たりしても、同じ答えになります。

\bkviihead{2乗の計算}
$(1+2j)^2$ は、3-8の $(a+b)^2=a^2+2ab+b^2$ を使います。
\bkviiline{$(1+2j)^2=1^2+2\times1\times2j+(2j)^2=1+4j+4j^2=1+4j-4=-3+4j$}
$(2j)^2=2^2\times j^2=4\times(-1)=-4$ です。まん中の項 $2ab$ を忘れないようにしましょう。

\bksubsection{7}{共役な複素数：かけると実数になる組}

\leavevmode\label{bk:6-8:kyoyaku}$3+4j$ と $3-4j$ のように、実部が同じで、虚部の符号だけが反対の2つの複素数を、たがいに\textbf{共役}（きょうやく）な複素数といいます。\bkidx{きょうやくふくそすう}{共役な複素数}もとは「くびき」（2頭の牛をならべてつなぐ横木）を表す字で書き、組になってはたらく2つ、という意味です。

\noindent\begin{minipage}[c]{0.6\linewidth}
\raggedright
平面で見ると、共役な2つの点は、横の線（実数の線）について線対称です（4-1）。横の線を鏡にして、うつした位置です。

$z$ と共役な複素数を、$\bar z$（ゼット・バー）と書きます。$z=a+bj$ なら $\bar z=a-bj$ です。$\bar z$ と共役な複素数は、もう一度裏返すので、もとの $z$ にもどります。
\end{minipage}\hfill
\begin{minipage}[c]{0.36\linewidth}
\centering
\begin{tikzpicture}[x=5.6mm, y=5.6mm]
  \draw[bkLine, line width=0.4pt, step=1] (-0.6,-4.6) grid (4.6,4.6);
  \bkviiaxes{-0.8}{4.9}{-4.8}{4.9}
  \draw[bkGray, line width=0.6pt, dash pattern=on 2pt off 1.5pt] (3,-4) -- (3,4);
  \bkviiarrow{bkBlue}{(3,4)}
  \bkviiarrow{bkOrange}{(3,-4)}
  \node[font=\sffamily\footnotesize, anchor=south, fill=white, inner sep=1pt] at (3,4.2) {$z=3+4j$};
  \node[font=\sffamily\footnotesize, anchor=north, fill=white, inner sep=1pt] at (3,-4.2) {$\bar z=3-4j$};
\end{tikzpicture}
\end{minipage}

\par\vspace{4pt}
共役な2つをかけてみます。3-8の「和と差の積」$(A+B)(A-B)=A^2-B^2$ で、$A=3$、$B=4j$ とします。
\begin{bkhenkei}
(3+4j)(3-4j)=3^2-(4j)^2 & 和と差の積（3-8） \\
=9-16j^2 & $(4j)^2=4^2\times j^2=16j^2$ \\
=9-16\times(-1) & $j^2=-1$ \\
=9+16=25 & \\
\end{bkhenkei}
$j$ が消えて、実数になりました。文字で書くと、
\bkviiframe[90mm]{$(a+bj)(a-bj)=a^2-b^2j^2=a^2+b^2$}
$a^2+b^2$ は、2乗の和なので、いつも0以上の実数です。この「かけると実数になる」性質を、次のわり算で使います。

\bksubsection{8}{わり算：分母の共役をかける}

\leavevmode\label{bk:6-8:waru}$\dfrac{5}{1+2j}$ は、どんな複素数でしょう。このままでは、分母に $j$ があって、$a+bj$ の形になっていません。

3-10では、$\dfrac{1}{\sqrt2}$ の分母の $\sqrt{\phantom{2}}$ をなくすために、分母と分子に同じ $\sqrt2$ をかけました（有理化）。分母と分子に同じ数をかけても、分数の大きさは変わらない（2-4）からです。ここでも同じ考えで、分母と分子に、\textbf{分母と共役な数} $1-2j$ をかけます。⑦の性質で、分母が実数になるからです。
\begin{bkhenkei}
\frac{5}{1+2j}=\frac{5(1-2j)}{(1+2j)(1-2j)} & 分母と分子に $1-2j$ をかける \\
=\frac{5(1-2j)}{1^2+2^2} & 分母は⑦の $a^2+b^2$ \\
=\frac{5(1-2j)}{5} & $1+4=5$ \\
=1-2j & 5で約分する \\
\end{bkhenkei}
\bkviihead{かけ算で確かめる}
わり算はかけ算の逆（1-14）なので、答えに分母をかけて、分子にもどるかを見ます。
\bkviiline{$(1-2j)(1+2j)=1^2+2^2=5$}
たしかに5にもどりました。

\bkviihead{$\dfrac{1}{j}$ を3通りで求める}
\begin{enumerate}[label={\bkcircled[\bkTF]{\arabic*}}, leftmargin=2.4\zw, itemsep=3pt, topsep=2pt]
\item \textbf{共役をかける}\quad $j=0+1j$ と共役な数は $-j$ です。分母と分子に $-j$ をかけると、
\bkviiline{$\dfrac{1}{j}=\dfrac{1\times(-j)}{j\times(-j)}=\dfrac{-j}{-j^2}=\dfrac{-j}{-(-1)}=\dfrac{-j}{1}=-j$}
\item \textbf{かけて1になる数をさがす}\quad $\dfrac{1}{j}$ は、$j$ にかけると1になる数です（逆数。2-4）。$j\times(-j)=-j^2=1$ なので、その数は $-j$ です。
\item \textbf{回転で考える}\quad $\times j$ は左回りに90°回すことでした。$\div j$ は、その逆なので、右回りに90°回すことです。1を右回りに90°回すと、真下へ1の所、つまり $-j$ です。
\end{enumerate}
3通りとも、$\dfrac{1}{j}=-j$ になりました。

\bksubsection{9}{2次方程式の、虚数の解}

\leavevmode\label{bk:6-8:kai}いよいよ、3-10で「解はありません」としていた方程式を解きます。

\bkviihead{$x^2=-4$}
$x^2=-4$ は、「2乗すると $-4$ になる数」です。$(2j)^2=2^2\times j^2=4\times(-1)=-4$ なので、$2j$ が解です。$(-2j)^2=(-2)^2\times j^2=4\times(-1)=-4$ なので、$-2j$ も解です。
\bkviiframe[60mm]{\large $x^2=-4$\quad の解は\quad $x=\pm2j$}
2乗して4になる数が $\pm2$ なのと同じで、2つの解は「プラス・マイナス」の組です。そこで、$a$ が正の数のとき、$\sqrt{-a}$ を次のように決めます。
\bkviiframe[80mm]{$a>0$ のとき\qquad $\sqrt{-a}=\sqrt a\,j$}
たとえば $\sqrt{-4}=2j$、$\sqrt{-16}=4j$、$\sqrt{-5}=\sqrt5\,j$ です。$x^2=-a$ の解は $x=\pm\sqrt a\,j$ です。

\bkviihead{$s^2+2s+5=0$（3-10の例）}
3-10の物理・工学ではこう使う（車のばねとダンパー）で出てきた2次方程式です。6-12では、ここで文字 $x$ のかわりに $s$ を使います（3-10と同じ）。2通りで解きます。

\noindent\textbf{方法1：解の公式で}\par
$a=1$、$b=2$、$c=5$ です。判別式は $D=2^2-4\times1\times5=4-20=-16$ で、マイナスです。
\begin{bkhenkei}
s=\frac{-2\pm\sqrt{-16}}{2\times1} & 解の公式（3-10） \\
=\frac{-2\pm4j}{2} & $\sqrt{-16}=\sqrt{16}\,j=4j$ \\
=-1\pm2j & 分子の両方の項を2でわる \\
\end{bkhenkei}
\noindent\textbf{方法2：平方完成で}\par
$s^2+2s=(s+1)^2-1$（3-8）なので、$s^2+2s+5=(s+1)^2+4$。だから、方程式は
\bkviiline{$(s+1)^2+4=0$\qquad つまり\qquad $(s+1)^2=-4$}
かっこの中身 $s+1$ は「2乗すると $-4$ になる数」なので、$s+1=\pm2j$。両辺から1をひいて $s=-1\pm2j$。方法1と同じ答えです。

\bkviihead{代入して確かめる}
$s=-1+2j$ を、もとの式の左辺に入れます。
\begin{bkhenkei}
s^2=(-1+2j)^2=1-4j+4j^2 & $(a+b)^2$ の展開（3-8） \\
=1-4j-4=-3-4j & $j^2=-1$ \\
2s=2(-1+2j)=-2+4j & 2をかける \\
s^2+2s+5 & 3つをたす \\
=(-3-4j)+(-2+4j)+5 & \\
=(-3-2+5)+(-4+4)j=0 & 実部どうし、虚部どうし \\
\end{bkhenkei}
たしかに0になりました。$s=-1-2j$ でも同じように0になります。ノートで確かめてみましょう。

\bkviihead{2つの解から、もとの式を作る}
3-10では、$ax^2+bx+c=0$ の解が $\alpha$、$\beta$ なら、左辺は $a(x-\alpha)(x-\beta)$ と因数分解できました。$\alpha=-1+2j$、$\beta=-1-2j$ で確かめます。$s-\alpha=(s+1)-2j$、$s-\beta=(s+1)+2j$ なので、和と差の積（3-8）で、
\bkviiline{$(s-\alpha)(s-\beta)=(s+1)^2-(2j)^2=s^2+2s+1+4=s^2+2s+5$}
もとの式にもどりました。展開すると $(s-\alpha)(s-\beta)=s^2-(\alpha+\beta)s+\alpha\beta$ なので、2つの解の和 $\alpha+\beta=-2$ と積 $\alpha\beta=(-1)^2+2^2=5$（⑦）が、式の係数の $-2$（$s$ の係数2の符号を変えたもの）と5に合っています。

\bkviihead{判別式と解の種類}
$D<0$ のとき、解の公式の $\sqrt{D}$ が純虚数 $\sqrt{-D}\,j$ になります。だから、2つの解は $p+qj$ と $p-qj$ の形の、共役な組になります（$p=-\dfrac{b}{2a}$、$q=\dfrac{\sqrt{-D}}{2a}$）。3-10の表を、複素数まで広げて書きなおします。
\begin{bkcenter}
\begin{bktable}{colspec={Q[c,m,wd=18mm] Q[l,m,wd=46mm] Q[l,m,wd=46mm]}, rowsep=1.5pt}
判別式 & 3-10（実数の中で） & 6-8（複素数まで広げて） \\
$D>0$ & 解は2つ & 実数の解が2つ \\
$D=0$ & 解は1つ & 実数の解が1つ（重なった解） \\
$D<0$ & 解はない & 虚数の解が2つ（共役な組） \\
\end{bktable}
\bktabcaption{3-10の「解はない」は、正しくは「実数の解はない」だった。}
\end{bkcenter}
6-12では、この表の3つの場合が、ばねにつるしたおもりや車の「ゆれずにもどる」「ちょうどよくもどる」「ゆれながらもどる」の3つに分かれることを学びます（\bkreidairef{8}）。

\bkviihead{グラフと平面で見る}
2つの式 $s^2+3s+2=0$ と $s^2+2s+5=0$ を、グラフと平面でくらべます（図\ref{bk:6-8:fig-kai}）。方法2の平方完成 $s^2+2s+5=(s+1)^2+4$ から、グラフ $y=s^2+2s+5$ の頂点は $(-1,\ 4)$ です（3-9）。いちばん低い所でも高さが4なので、横軸とは交わりません。ここで、解 $-1\pm2j$ をよく見ると、実部の $-1$ は頂点の横の位置、虚部の2は、2乗すると頂点の高さ4になる数です。$(s+1)^2=-4$ の「4」が、そのまま頂点の高さになっているからです。

\begin{bkfigure}
\begin{tikzpicture}[baseline=(current bounding box.center)]
\begin{axis}[bkplot, name=G, width=64mm, height=52mm, xmin=-4.2, xmax=2.2, ymin=-1.5, ymax=10.5, xtick={-3,-2,-1,1,2}, ytick={4,8}, xlabel={$s$}, ylabel={$y$},
    yticklabel style={fill=white, inner sep=1pt, outer sep=2.3pt}]
  \addplot[bkcurve, domain=-3.5:1.6, samples=60] {x^2+2*x+5};
  \addplot[bkcurve2, domain=-4.1:1.1, samples=60] {x^2+3*x+2};
  \addplot[only marks, mark=*, mark size=1.8pt, bkOrange] coordinates {(-1,0) (-2,0)};
  \node[font=\sffamily\footnotesize, text=bkBlue, fill=white, inner sep=1pt] at (axis cs:-1.4,9.3) {$s^2+2s+5$};
  \node[font=\sffamily\footnotesize, text=bkOrange!80!black, anchor=east, fill=white, inner sep=1pt] at (axis cs:-3.75,6.0) {$s^2+3s+2$};
\end{axis}
\end{tikzpicture}\hspace{8mm}%
\begin{tikzpicture}[x=9mm, y=9mm, baseline=(current bounding box.center)]
  \draw[bkLine, line width=0.4pt, step=1] (-3.5,-2.5) grid (1.5,2.5);
  \bkviiaxes{-3.7}{1.8}{-2.7}{2.8}
  \foreach \k in {-3,-2,-1,1} {\node[font=\sffamily\footnotesize, anchor=north, inner sep=1pt, fill=white] at (\k,-0.12) {$\k$};}
  \node[font=\sffamily\footnotesize, anchor=west, inner sep=1pt, fill=white] at (0.1,2) {$2j$};
  \node[font=\sffamily\footnotesize, anchor=west, inner sep=1pt, fill=white] at (0.1,-2) {$-2j$};
  \fill[bkOrange] (-1,0) circle (0.1) (-2,0) circle (0.1);
  \fill[bkBlue] (-1,2) circle (0.1) (-1,-2) circle (0.1);
  \node[font=\sffamily\footnotesize, anchor=south east, fill=white, inner sep=1pt] at (-1.1,2.1) {$-1+2j$};
  \node[font=\sffamily\footnotesize, anchor=north east, fill=white, inner sep=1pt] at (-1.1,-2.1) {$-1-2j$};
\end{tikzpicture}
\bkfigcaption{左：$y=s^2+3s+2$（オレンジ）は横軸と $-2$、$-1$ で交わる。$y=s^2+2s+5$（青）は交わらない。右：解を平面にかく。$D>0$ の解（オレンジ）は横の線の上、$D<0$ の解（青）は線の外に、上下対称にある。}\label{bk:6-8:fig-kai}
\end{bkfigure}

$y=s^2+2s+5$ のグラフが横軸と交わらないのは、解が「横の線（実数）の上にない」からでした。解は消えたのではなく、平面の上の、横の線からはなれた所にあったのです。6-12では、この $-1$ と $2$ が、車のゆれの「止まっていく速さ」と「ゆれの速さ」を表すことがわかります（物理・工学ではこう使う）。

\bksubsection{10}{コンピュータで複素数を計算する}

\leavevmode\label{bk:6-8:pc}Python では、虚数単位を \texttt{j} と書き、数のすぐあとにつけます。$j$ そのものは \texttt{1j} と書きます（\texttt{j} だけだと、文字の名前と区別がつかないため）。

\noindent\begin{minipage}[t]{0.5\linewidth}\vspace{0pt}
\begin{bkcode}[Python]
import cmath
z = 3+4j
w = 1-2j
print(z + w)
print(z * w)
print(z / w)
print(z.real, z.imag)
print(z.conjugate())
print(1j**2)
print(cmath.sqrt(-4))
\end{bkcode}
\end{minipage}\hfill
\begin{minipage}[t]{0.46\linewidth}\vspace{0pt}
\begin{bkcodeout}
(4+2j)
(11-2j)
(-1+2j)
3.0 4.0
(3-4j)
(-1+0j)
2j
\end{bkcodeout}
\end{minipage}

\par\vspace{4pt}
\texttt{z.real} は実部、\texttt{z.imag} は虚部、\texttt{z.conjugate()} は共役な複素数です。\texttt{**} は累乗、\texttt{cmath.sqrt} は複素数まで考えた $\sqrt{\phantom{2}}$ です。⑥のかけ算の答え $11-2j$ が、そのまま出てきました。\texttt{z / w} の答え $-1+2j$ が正しいかどうかは、\bkreidairef{5}で手で確かめます。

表計算ソフトにも、複素数を文字の形 \texttt{"3+4j"} で計算する関数があります。
\begin{bkcenter}
\begin{bktable}{colspec={Q[l,m,wd=50mm] Q[l,m,wd=30mm] Q[l,m,wd=30mm]}, rowsep=1pt, cells={font=\small}}
入れる式 & 答え & 意味 \\
\texttt{=COMPLEX(3,4,"j")} & \texttt{3+4j} & 複素数を作る \\
\texttt{=IMSUM("3+4j","1-2j")} & \texttt{4+2j} & たし算 \\
\texttt{=IMPRODUCT("3+4j","1-2j")} & \texttt{11-2j} & かけ算 \\
\texttt{=IMDIV("5","1+2j")} & \texttt{1-2j} & わり算 \\
\texttt{=IMREAL("3+4j")} & \texttt{3} & 実部 \\
\texttt{=IMAGINARY("3+4j")} & \texttt{4} & 虚部 \\
\end{bktable}
\bktabcaption{表計算の複素数の関数（\texttt{"j"} のかわりに \texttt{"i"} でもよい）}
\end{bkcenter}

\begin{motto}[4]
\textbf{複素数には、大きい・小さいがない}\par
$j$ は0より大きいでしょうか、小さいでしょうか。どちらと決めても、おかしなことが起きます。
\begin{itemize}[leftmargin=1.2\zw, itemsep=2pt, topsep=2pt, label={\color{bkGray}●}]
\item $j>0$ と決めると：正の数どうしをかけると正なので、$j\times j>0$。ところが $j\times j=-1$ で、$-1>0$ になってしまう。
\item $j<0$ と決めると：$-j>0$ になるので、$(-j)\times(-j)>0$。ところが $(-j)^2=-1$ で、やはり $-1>0$ になってしまう。
\end{itemize}
だから、虚数には「大きい・小さい」を決めません。$3+4j>2+j$ のような式は、意味がありません。複素数の「大きさ」をくらべたいときは、原点からの距離（絶対値。6-9）をくらべます。
\end{motto}

\begin{motto}[4]
\textbf{複素数まで広げれば、方程式の解はいつもある}\par
実数だけで考えると、2次方程式には解が2つ・1つ・0の場合がありました。複素数まで広げると、2次方程式の解は、いつでも2つです（$D=0$ のときは、同じ解が2つ重なったと数えます）。3次方程式なら3つ、$n$ 次方程式なら $n$ この解が、複素数の中にかならずあります。これを\textbf{代数学の基本定理}といいます。ガウスが1799年の学位論文で証明を示し、そのあと多くの人が、よりきちんとした証明を考えました。
\par\vspace{3pt}
「2乗して $-1$」のために数を広げたら、もうそれ以上、方程式のために新しい数を考える必要はなくなったのです。第9巻の制御では、3次の特性方程式の解をさがす場面が出てきます。6-9では、$z^3=1$ の3つの解が、平面の上で正三角形の頂点にならぶことを見ます。
\end{motto}




\begin{kotoba}
\small
\bkword{虚数単位}{2乗すると $-1$ になる数 $j$。かけると、左回りに90°回る。}
\bkword{複素数}{$a+bj$（$a$、$b$ は実数）の形の数。平面の点 $(a,\ b)$ で表せる。}
\bkword{実部・虚部}{$a+bj$ の $a$ と $b$。虚部は $bj$ ではなく $b$。$z$ の実部を $\mathrm{Re}\,z$、虚部を $\mathrm{Im}\,z$ と書く。}
\bkword{虚数・純虚数}{虚部が0でない複素数が虚数。そのうち実部が0のものが純虚数。}
\bkword{共役な複素数}{$a+bj$ と $a-bj$ のように、虚部の符号だけが反対の組。$\bar z$ と書く。}
\end{kotoba}

\begin{senmonsho}[大学の教科書ふうの書き方]
\begin{bkhonbun}
$i^2=-1$ をみたす数 $i$ を虚数単位という。実数 $x$、$y$ を用いて $z=x+iy$ と表される数を複素数といい、$x=\mathrm{Re}\,z$ を実部、$y=\mathrm{Im}\,z$ を虚部という。$\bar z=x-iy$ を $z$ の共役複素数という。$z\bar z=x^2+y^2$ は非負の実数である。
\end{bkhonbun}
\bkyaku{数学の本では、虚数単位を $i$ と書き、$x+iy$ のように $i$ を前に書くことも多い。「用いて」は「使って」。「非負」は「マイナスでない」（0以上）。$z$ に $\bar z$ をかけると、0以上の実数になる。}
\end{senmonsho}

\bkkigocards{\bkkigocard{$j$}{ジェイ}{虚数単位。2乗すると $-1$。数学の本・高校では $i$ と書く}{imaginary unit}{6-8}}
{\bkkigocard{$\bar z$}{ゼット・バー}{$z$ と共役な複素数。$z=a+bj$ なら $\bar z=a-bj$}{complex conjugate}{6-8}}

\begin{taisetsu}
\def\tsl#1{\makebox[6.2\zw][l]{\sffamily\bfseries\small\color{\bkTT}#1}}
\tsl{虚数単位}$j^2=-1$。$\times j$ は、左回りに90°回すこと\\[3pt]
\tsl{複素数}$a+bj$（$a$ は実部、$b$ は虚部）。平面の点 $(a,\ b)$\\[3pt]
\tsl{計算}$j$ を文字のように計算し、$j^2$ が出たら $-1$ にする\\[3pt]
\tsl{共役}$(a+bj)(a-bj)=a^2+b^2$。わり算は、分母と共役な数を分母・分子にかける\\[3pt]
\tsl{2次方程式}$D<0$ なら、虚数の解が2つ（共役な組）。$\sqrt{-a}=\sqrt a\,j$
\end{taisetsu}



%% ---------------- 例題 ----------------
\bksection{reidai}{例題でたしかめよう}

\begin{reidai}[$j$ の累乗]
次の数を、$1$、$j$、$-1$、$-j$ のどれかで答えましょう。\newline\bkviiit{（1）$j^3$}\bkviiit{（2）$j^8$}\bkviiit{（3）$j^{10}$}\bkviiit{（4）$j^{13}$}\par
\kangaekata $j$ の累乗は、4回でもとにもどる（④）。指数を4でわったあまりを見る。あまりが1なら $j$、2なら $-1$、3なら $-j$、0なら1。\par
\shiki\hangindent=4.3\zw\hangafter=1 （1）$j^3=j^2\times j=(-1)\times j=-j$\par
\bkviiq （2）$8=4\times2$ あまり0。$j^8=(j^4)^2=1^2=1$\par
\bkviiq （3）$10=4\times2+2$ あまり2。$j^{10}=(j^4)^2\times j^2=1\times(-1)=-1$\par
\bkviiq （4）$13=4\times3+1$ あまり1。$j^{13}=(j^4)^3\times j=1\times j=j$\par
\kotae\hangindent=4.3\zw\hangafter=1 （1）$-j$\quad（2）1\quad（3）$-1$\quad（4）$j$
\end{reidai}

\begin{reidai}[たし算・ひき算・実数倍]
$z=3+4j$、$w=1-2j$ のとき、次の計算をしましょう。\newline\bkviiit{（1）$z+w$}\bkviiit{（2）$z-w$}\bkviiit{（3）$2z-3w$}\par
\kangaekata 実部どうし、虚部どうしをまとめる（⑤）。かっこの前の $-$ は、かっこの中の両方の項につく。\par
\shiki\hangindent=4.3\zw\hangafter=1 （1）$(3+4j)+(1-2j)=(3+1)+(4-2)j=4+2j$\par
\bkviiq （2）$(3+4j)-(1-2j)=3+4j-1+2j=2+6j$\par
\bkviiq （3）$2z-3w=(6+8j)-(3-6j)=(6-3)+(8+6)j=3+14j$\par
\kotae\hangindent=4.3\zw\hangafter=1 （1）$4+2j$\quad（2）$2+6j$\quad（3）$3+14j$
\end{reidai}

\noindent\begin{minipage}[t]{0.485\linewidth}\vspace{0pt}
\begin{jibun}
次の数を、$1$、$j$、$-1$、$-j$ のどれかで答えましょう。\par
（1）$j^5$\newline（2）$j^7$\newline（3）$j^{12}$\newline（4）$j^{22}$
\bkjans{（1）\nobreak $5=4+1$ で $j$。（2）\nobreak $7=4+3$ で $-j$。（3）\nobreak $12=4\times3$ で1。（4）\nobreak $22=4\times5+2$ で $-1$。}
\end{jibun}
\end{minipage}\hfill
\begin{minipage}[t]{0.485\linewidth}\vspace{0pt}
\begin{jibun}
$z=2-5j$、$w=-1+3j$ のとき、次の計算をしましょう。\par
（1）$z+w$\newline（2）$z-w$\newline（3）$3z+2w$
\bkjans{（1）\nobreak $1-2j$。（2）\nobreak $3-8j$。（3）\nobreak $(6-15j)+(-2+6j)=4-9j$。}
\end{jibun}
\end{minipage}

\begin{reidai}[かけ算]
次の計算をしましょう。\newline\bkviiit[0.32]{（1）$(3+4j)(1-2j)$}\bkviiit[0.32]{（2）$(1+2j)^2$}\bkviiit[0.32]{（3）$j(2+5j)$}\par
\kangaekata 展開して、$j^2$ を $-1$ にする（⑥）。\par
（2）は $(a+b)^2=a^2+2ab+b^2$ を使う（3-8）。\par
\shiki\hangindent=4.3\zw\hangafter=1 （1）$3-6j+4j-8j^2=3-2j-8\times(-1)=3+8-2j$\par
\bkviiq （2）$1^2+2\times1\times2j+(2j)^2=1+4j+4j^2=1+4j-4$\par
\bkviiq （3）$2j+5j^2=2j-5$。確かめ：点 $(2,\ 5)$ を90°回すと $(-5,\ 2)$（4-8）\par
\kotae\hangindent=4.3\zw\hangafter=1 （1）$11-2j$\quad（2）$-3+4j$\quad（3）$-5+2j$
\end{reidai}

\begin{reidai}[共役な複素数をかける]
次の複素数 $z$ について、共役な複素数 $\bar z$ と、積 $z\bar z$ を求めましょう。\newline\bkviiit{（1）$z=2+5j$}\bkviiit{（2）$z=-3-j$}\par
\kangaekata 虚部の符号を反対にすると $\bar z$。\par
積は $(a+bj)(a-bj)=a^2+b^2$（⑦）。\par
\shiki\hangindent=4.3\zw\hangafter=1 （1）$\bar z=2-5j$。$z\bar z=2^2+5^2=4+25=29$\par
\bkviiq （2）$\bar z=-3+j$。$a=-3$、$b=-1$ なので $z\bar z=(-3)^2+(-1)^2=10$\par
\bkviib 確かめ：$(-3-j)(-3+j)=9-3j+3j-j^2=9+1=10$\par
\kotae\hangindent=4.3\zw\hangafter=1 （1）$\bar z=2-5j$、$z\bar z=29$\quad（2）$\bar z=-3+j$、$z\bar z=10$
\end{reidai}

\noindent\begin{minipage}[t]{0.485\linewidth}\vspace{0pt}
\begin{jibun}
次の計算をしましょう。\par
（1）$(2+3j)(1+j)$\newline（2）$(3-j)^2$\par
（3）$j(4-3j)$
\bkjans{（1）\nobreak $2+2j+3j+3j^2=-1+5j$。（2）\nobreak $9-6j+j^2=8-6j$。（3）\nobreak $4j-3j^2=3+4j$。}
\end{jibun}
\end{minipage}\hfill
\begin{minipage}[t]{0.485\linewidth}\vspace{0pt}
\begin{jibun}
次の $z$ の $\bar z$ と $z\bar z$ を求めましょう。\par
（1）$z=4-3j$\newline（2）$z=-1+2j$
\bkjans{（1）\nobreak $\bar z=4+3j$、$z\bar z=16+9=25$。（2）\nobreak $\bar z=-1-2j$、$z\bar z=1+4=5$。}
\end{jibun}
\end{minipage}

\noindent かけ算の答えは、⑩の Python で \texttt{(2+3j)*(1+j)} のように計算して確かめられます。$z\bar z$ は、いつも「実部の2乗＋虚部の2乗」の、0以上の実数になります。答えに $j$ が残ったり、マイナスになったりしたら、どこかで $j^2=-1$ の使い方をまちがえています。

\begin{reidai}[わり算]
次の計算をしましょう。\newline\bkviiit{（1）$\dfrac{5}{1+2j}$}\bkviiit{（2）$\dfrac{3+4j}{1-2j}$}\bkviiit{（3）$\dfrac{2}{j}$}\par
\kangaekata 分母と共役な数を、分母と分子にかける（⑧）。分母は $a^2+b^2$ の実数になる。\par
\shiki\hangindent=4.3\zw\hangafter=1 （1）$\dfrac{5(1-2j)}{(1+2j)(1-2j)}=\dfrac{5(1-2j)}{1+4}=1-2j$\par
\bkviiq （2）分母と分子に $1+2j$ をかける。分母は $1^2+2^2=5$\par
\bkviib 分子は $(3+4j)(1+2j)=3+6j+4j+8j^2=-5+10j$\par
\bkviib $\dfrac{-5+10j}{5}=-1+2j$。確かめ：$(-1+2j)(1-2j)=3+4j$\par
\bkviiq （3）分母と分子に $-j$ をかける。$\dfrac{2\times(-j)}{j\times(-j)}=\dfrac{-2j}{-j^2}=-2j$\par
\kotae\hangindent=4.3\zw\hangafter=1 （1）$1-2j$\quad（2）$-1+2j$\quad（3）$-2j$
\end{reidai}

\begin{reidai}[$x^2=-a$ の形と、マイナスの数の $\sqrt{\phantom{2}}$]
（1）$x^2+9=0$ を解きましょう。\newline（2）$x^2=-5$ を解きましょう。\newline（3）$\sqrt{-4}\times\sqrt{-9}$ を計算しましょう。\par
\kangaekata $a>0$ のとき $\sqrt{-a}=\sqrt a\,j$（⑨）。マイナスの数の $\sqrt{\phantom{2}}$ は、先に $j$ を使った形に直す。\par
\shiki\hangindent=4.3\zw\hangafter=1 （1）$x^2=-9$ なので $x=\pm\sqrt9\,j=\pm3j$（確かめ：$(3j)^2=9j^2=-9$）\par
\bkviiq （2）$x=\pm\sqrt5\,j$\par
\bkviiq （3）$\sqrt{-4}\times\sqrt{-9}=2j\times3j=6j^2=-6$\par
\kotae\hangindent=4.3\zw\hangafter=1 （1）$x=\pm3j$\quad（2）$x=\pm\sqrt5\,j$\quad（3）$-6$
\end{reidai}

\noindent\begin{minipage}[t]{0.485\linewidth}\vspace{0pt}
\begin{jibun}
次の計算をしましょう。\par
（1）$\dfrac{10}{3-j}$\newline（2）$\dfrac{1+3j}{1+j}$\newline（3）$\dfrac{3}{j}$
\bkjans{（1）\nobreak $\dfrac{10(3+j)}{9+1}=3+j$。（2）\nobreak $\dfrac{(1+3j)(1-j)}{1+1}=\dfrac{4+2j}{2}=2+j$。（3）\nobreak $\dfrac{3\times(-j)}{j\times(-j)}=-3j$。}
\end{jibun}
\end{minipage}\hfill
\begin{minipage}[t]{0.485\linewidth}\vspace{0pt}
\begin{jibun}
（1）（2）を解き、（3）を計算しましょう。\newline（1）$x^2+16=0$\newline（2）$x^2=-7$\newline（3）$\sqrt{-2}\times\sqrt{-8}$
\bkjans{（1）\nobreak $x=\pm4j$。（2）\nobreak $x=\pm\sqrt7\,j$。（3）\nobreak $\sqrt2\,j\times\sqrt8\,j=\sqrt{16}\,j^2=-4$。}
\end{jibun}
\end{minipage}

\begin{reidai}[解の公式で、虚数の解を求める]
$s^2+4s+13=0$ を解きましょう。また、2つの解の和と積を求めて、式の係数と合うか確かめましょう。\par
\kangaekata 解の公式（3-10）と $\sqrt{-a}=\sqrt a\,j$。2つの解を $\alpha$、$\beta$ とすると $s^2+bs+c=(s-\alpha)(s-\beta)$ なので、和は $-b$、積は $c$（⑨）。\par
\shiki\hangindent=2.6\zw\hangafter=1 $a=1$、$b=4$、$c=13$。$D=4^2-4\times1\times13=16-52=-36$\par
\bkviia $s=\dfrac{-4\pm\sqrt{-36}}{2}=\dfrac{-4\pm6j}{2}=-2\pm3j$\par
\bkviia 和は $-4$（係数4の符号を変えた数）。積は $(-2)^2+3^2=13$（数の項）\par
\kotae\hangindent=2.6\zw\hangafter=1 $s=-2\pm3j$（和 $-4$、積13で、合っている）
\end{reidai}

\begin{reidai}[判別式で、解の種類を見分ける]
次の2次方程式の判別式 $D$ を求め、解の種類を答えてから、解きましょう。\par
\noindent \bkviiit[0.32]{（1）$s^2+3s+2=0$}\bkviiit[0.32]{（2）$s^2+4s+4=0$}\bkviiit[0.32]{（3）$s^2+s+1=0$}\par
\kangaekata $D=b^2-4ac$ の符号で、解の種類がわかる（⑨の表）。\par
\shiki\hangindent=4.3\zw\hangafter=1 （1）$D=9-8=1>0$。$(s+1)(s+2)=0$ より $s=-1$、$-2$\par
\bkviiq （2）$D=16-16=0$。$(s+2)^2=0$ より $s=-2$\par
\bkviiq （3）$D=1-4=-3<0$。$s=\dfrac{-1\pm\sqrt{-3}}{2}=\dfrac{-1\pm\sqrt3\,j}{2}$\par
\kotae\hangindent=4.3\zw\hangafter=1 （1）実数の解2つ。$s=-1$、$-2$\quad（2）実数の解1つ。$s=-2$\par
\bkviiq （3）虚数の解2つ。$s=-\dfrac{1}{2}\pm\dfrac{\sqrt3}{2}j$
\end{reidai}

\noindent\begin{minipage}[t]{0.485\linewidth}\vspace{0pt}
\begin{jibun}
$s^2-4s+13=0$ を解きましょう。また、2つの解の和と積を求めて、式の係数と合うか確かめましょう。
\bkjans{$D=16-52=-36$。$s=\dfrac{4\pm6j}{2}=2\pm3j$。和4、積 $4+9=13$ で、合っている。}
\end{jibun}
\end{minipage}\hfill
\begin{minipage}[t]{0.485\linewidth}\vspace{0pt}
\begin{jibun}
判別式 $D$ を求め、解の種類を答えてから、解きましょう。\par
（1）$s^2+5s+6=0$\newline（2）$s^2+6s+9=0$\par
（3）$s^2+2s+2=0$
\bkjans{（1）\nobreak $D=1>0$ で実数の解2つ。$s=-2$、$-3$。（2）\nobreak $D=0$ で実数の解1つ。$s=-3$。（3）\nobreak $D=-4<0$ で虚数の解2つ。$s=\dfrac{-2\pm2j}{2}=-1\pm j$。}
\end{jibun}
\end{minipage}

\begin{reidai}[和と積から、2つの数を求める]
「10を2つに分けて、かけると40になるようにせよ」（コラムの、カルダーノの問題）。2つの数を求め、たすと10、かけると40になることを確かめましょう。\par
\kangaekata 1つの数を $x$ とすると、もう1つは $10-x$。$x(10-x)=40$ を2次方程式の形にして、解の公式（3-10）で解く。\par
\shiki\hangindent=2.6\zw\hangafter=1 左辺を展開して $10x-x^2=40$。両辺から $10x-x^2$ をひいて $0=x^2-10x+40$\par
\bkviia $D=(-10)^2-4\times1\times40=100-160=-60$\par
\bkviia $\sqrt{-60}=\sqrt{60}\,j=2\sqrt{15}\,j$（$60=4\times15$、$\sqrt4=2$）\par
\bkviia $x=\dfrac{10\pm2\sqrt{15}\,j}{2}=5\pm\sqrt{15}\,j$。$x=5+\sqrt{15}\,j$ なら、もう1つは $10-x=5-\sqrt{15}\,j$\par
\bkviia 確かめ：和は $5+5=10$。積は $5^2+(\sqrt{15})^2=25+15=40$（⑦）\par
\kotae\hangindent=2.6\zw\hangafter=1 $5+\sqrt{15}\,j$ と $5-\sqrt{15}\,j$
\end{reidai}

\begin{jibun}
たすと6、かけると13になる2つの数を求めましょう。また、たすと6、かけると13になることを確かめましょう。
\bkjans{$x(6-x)=13$ より $x^2-6x+13=0$。$D=36-52=-16$、$x=\dfrac{6\pm4j}{2}=3\pm2j$。2つの数は $3+2j$ と $3-2j$。和6、積 $9+4=13$。}
\end{jibun}

%% ---------------- まちがえやすいところ・物理工学 ----------------
\bksection{tsumazuki}{まちがえやすいところをチェック}
\noindent\begin{minipage}[t]{0.485\linewidth}\vspace{0pt}
\begin{tsumazuki}{$\sqrt{\phantom{2}}$ の中どうしをかける}
\bkng{$\sqrt{-4}\sqrt{-9}=\sqrt{36}=6$}\bkok{$2j\times3j=6j^2=-6$}
$\sqrt a\sqrt b=\sqrt{ab}$ は、$a$、$b$ が正の数のときに作ったきまり（2-8）です。先に $j$ の形に直します。
\end{tsumazuki}
\end{minipage}\hfill
\begin{minipage}[t]{0.485\linewidth}\vspace{0pt}
\begin{tsumazuki}{$j^2$ の符号}
\bkng{$(3j)^2=9$}\bkok{$(3j)^2=9j^2=-9$}
$j^2$ は $+1$ ではなく $-1$ です。$j$ が残ったら、すぐ $j^2=-1$ を使います。
\end{tsumazuki}
\end{minipage}
\par\vspace{2pt}
\noindent\begin{minipage}[t]{0.485\linewidth}\vspace{0pt}
\begin{tsumazuki}{虚部はどれ？}
\bkng{$3+4j$ の虚部は $4j$}\bkok{$3+4j$ の虚部は4}
虚部は、$j$ にかかっている実数です。実部も虚部も、どちらも実数です。
\end{tsumazuki}
\end{minipage}\hfill
\begin{minipage}[t]{0.485\linewidth}\vspace{0pt}
\begin{tsumazuki}{わり算で、分子だけにかける}
\bkng{$\dfrac{5}{1+2j}=\dfrac{5(1-2j)}{1+2j}$}\bkok{$\dfrac{5(1-2j)}{(1+2j)(1-2j)}$}
分子だけにかけると、数の大きさが変わってしまいます。分母と分子の両方にかけます。
\end{tsumazuki}
\end{minipage}

\begin{butsuri}
\bkviibox{交流：90°のずれを $j$ で書く}
\noindent\begin{minipage}[c]{0.5\linewidth}
\raggedright
コンデンサーに流れる電流は、電圧より $\frac{1}{4}$ 周期（90°）先に山が来ます（5-15）。電気の計算では、この「90°進む」を「$j$ をかける」で表します。電圧と電流を、回る矢印の組として書くと、ずれのある交流のたし算やわり算が、複素数の計算でできるのです（6-10・6-13、物理シリーズ『電磁気学』）。
\end{minipage}\hfill
\begin{minipage}[c]{0.46\linewidth}
\centering
\begin{tikzpicture}[x=1mm, y=1mm]
  \draw[bkGray, line width=0.6pt, -{Stealth[length=3pt]}] (-3,0) -- (24,0);
  \draw[bkGray, line width=0.6pt, -{Stealth[length=3pt]}] (0,-3) -- (0,22);
  \draw[bkBlue, line width=1.5pt, -{Stealth[length=4.5pt]}] (0,0) -- (19,0);
  \draw[bkOrange, line width=1.5pt, -{Stealth[length=4.5pt]}] (0,0) -- (0,16);
  \draw[bkGray, line width=0.7pt, dash pattern=on 2pt off 1.5pt, -{Stealth[length=3.5pt]}] (8:9) arc (8:82:9);
  \node[font=\sffamily\footnotesize, anchor=north, text=bkBlue] at (12,-1.6) {電圧};
  \node[font=\sffamily\footnotesize, anchor=west, text=bkOrange!80!black] at (1.6,18) {電流（$\times j$）};
  \node[font=\sffamily\footnotesize, text=bkGray] at (45:12.5) {90°};
\end{tikzpicture}
\end{minipage}
\end{butsuri}

\begin{butsuri}
\bkviibox{車のゆれ：解の実部と虚部が、ゆれ方を決める}
3-10の車のばねとダンパーの式 $s^2+2s+5=0$ の解は $s=-1\pm2j$ でした。6-12では、このとき車の高さのずれが $e^{-t}(\cos2t+\frac{1}{2}\sin2t)$ になることを学びます。実部の $-1$ が $e^{-t}$（だんだん小さくなる速さ）に、虚部の2が $\cos2t$（ゆれの速さ）に出てきます。解が実数の $s^2+3s+2=0$（$s=-1$、$-2$）では、ゆれずにもどります。
\begin{bkcenter}
\begin{tikzpicture}
\begin{axis}[bkplot, width=112mm, height=48mm, xmin=0, xmax=5.6, ymin=-0.6, ymax=1.15, xtick={1,2,3,4,5}, ytick={1}, xlabel={$t$（s）}, ylabel={高さのずれ},
    yticklabel style={fill=bkPurpleBg, inner sep=1pt, outer sep=2.3pt}]
  \addplot[bkcurve, domain=0:5.4, samples=150] {exp(-x)*(cos(deg(2*x))+0.5*sin(deg(2*x)))};
  \addplot[bkcurve2, domain=0:5.4, samples=80] {2*exp(-x)-exp(-2*x)};
  \node[font=\sffamily\footnotesize, text=bkBlue, anchor=west, fill=bkPurpleBg, inner sep=1pt] at (axis cs:1.75,-0.42) {$s=-1\pm2j$：ゆれながらもどる};
  \node[font=\sffamily\footnotesize, text=bkOrange!80!black, anchor=west, fill=bkPurpleBg, inner sep=1pt] at (axis cs:1.4,0.62) {$s=-1$、$-2$：ゆれずにもどる};
\end{axis}
\end{tikzpicture}
\end{bkcenter}
\end{butsuri}

\begin{butsuri}
\bkviibox{スマートフォンと Wi-Fi の電波}
\noindent\begin{minipage}[c]{0.6\linewidth}
\raggedright
スマートフォンや Wi-Fi の電波は、波の「大きさ」と「ずれ（位相）」の組で、情報を送っています。この組は、平面の点、つまり複素数で表せます。となりの点の図は、16この点を使い分ける方式（16QAM）で、1回に4ビット（$2^4=16$ 通り。2-17）を送ります。受け取る側は、届いた波から複素数を計算し、いちばん近い点を選んで、0と1にもどします。実際の Wi-Fi では、もっと多くの点を使う方式も使われています。
\end{minipage}\hfill
\begin{minipage}[c]{0.36\linewidth}
\centering
\begin{tikzpicture}[x=4.6mm, y=4.6mm]
  \draw[bkGray, line width=0.6pt, -{Stealth[length=3pt]}] (-4.4,0) -- (4.6,0);
  \draw[bkGray, line width=0.6pt, -{Stealth[length=3pt]}] (0,-4.4) -- (0,4.6);
  \foreach \x in {-3,-1,1,3} {\foreach \y in {-3,-1,1,3} {\fill[bkBlue] (\x,\y) circle (0.22);}}
  \node[font=\sffamily\footnotesize, anchor=north west, inner sep=1pt] at (4.2,-0.1) {実部};
  \node[font=\sffamily\footnotesize, anchor=south west, inner sep=1pt] at (0.15,4.3) {虚部};
\end{tikzpicture}
\end{minipage}
\end{butsuri}

\begin{column}{カルダーノの問題（1545年）}
イタリアのカルダーノは、3次方程式の解き方をまとめた本の中で、「10を2つに分けて、かけると40になるようにせよ」という問題を考えました。2つの数を $x$ と $10-x$ とすると、$x(10-x)=40$、つまり $x^2-10x+40=0$。解の公式で、$x=5\pm\sqrt{-15}$ が出てきます。カルダーノは、$(5+\sqrt{-15})(5-\sqrt{-15})=25-(-15)=40$ と確かめながらも、「役に立たない、頭をなやませるだけのもの」と書きました。実数の中では、たして10になる2つの数の積は、5と5のときの25がいちばん大きいので、40にはなりません。答えは、実数の外にしかなかったのです（\bkreidairef{9}）。
\end{column}

\begin{column}{「ないはずの数」が、本物の答えを連れてきた}
\textbf{ボンベリ（1572年）}\par
イタリアのボンベリは、$x^3=15x+4$ という3次方程式に、カルダーノの解き方を使いました。答えの1つは $x=4$ です（$4^3=64$、$15\times4+4=64$）。ところが計算のとちゅうで、$\sqrt{-121}=11j$ が出てきます。ボンベリは、虚数のまま計算を進めるきまりを作り、$(2+j)^3=2+11j$ を使って、虚数を通りぬけて本物の答え4にたどりつきました。
デカルト（1637年）がこうした数を「想像上の（imaginary）数」とよんだことが、「虚数」という名前のもとです。「ないはずの数」は、いまでは電気・制御・通信に欠かせない道具です。
\end{column}

%% ---------------- 章末問題 ----------------
\bkendproblems
\bkproblemgroup{基本}{4}

\bkproblem 次の数を、$1$、$j$、$-1$、$-j$ のどれかで答えましょう\mbox{（\bkreidairef{1}）}。\newline\bkviiit{（1）$j^6$}\bkviiit{（2）$j^9$}\bkviiit{（3）$j^{15}$}\bkviiit{（4）$j^{20}$}
\bkpans{（1）\nobreak $-1$。（2）\nobreak $j$。（3）\nobreak $-j$。（4）\nobreak 1。}

\bkproblem $z=5+2j$、$w=-2+j$ のとき、次の計算をしましょう\mbox{（\bkreidairef{2}）}。\newline\bkviiit{（1）$z+w$}\bkviiit{（2）$z-w$}\bkviiit{（3）$2z+3w$}
\bkpans{（1）\nobreak $3+3j$。（2）\nobreak $7+j$。（3）\nobreak $(10+4j)+(-6+3j)=4+7j$。}

\bkproblem 次の計算をしましょう\mbox{（\bkreidairef{3}）}。\newline\bkviiit[0.3]{（1）$(1+3j)(2-j)$}\bkviiit{（2）$(2+j)^2$}\bkviiit{（3）$j(1-4j)$}
\bkpans{（1）\nobreak $2-j+6j-3j^2=5+5j$。（2）\nobreak $4+4j+j^2=3+4j$。（3）\nobreak $j-4j^2=4+j$。}

\bkproblem $z=3-2j$ の共役な複素数 $\bar z$ と、積 $z\bar z$ を求めましょう\mbox{（\bkreidairef{4}）}。
\bkpans{$\bar z=3+2j$、$z\bar z=9+4=13$。}

\bkproblem 次の計算をしましょう\mbox{（\bkreidairef{5}）}。\newline\bkviiit{（1）$\dfrac{5}{2+j}$}\bkviiit{（2）$\dfrac{4+2j}{1+j}$}\bkviiit{（3）$\dfrac{4}{j}$}
\bkpans{（1）\nobreak $\dfrac{5(2-j)}{4+1}=2-j$。（2）\nobreak $\dfrac{(4+2j)(1-j)}{2}=\dfrac{6-2j}{2}=3-j$。（3）\nobreak $-4j$。}

\bkproblem 次の方程式を解きましょう\mbox{（\bkreidairef{6}）}。\newline\bkviiit{（1）$x^2+25=0$}\bkviiit{（2）$x^2=-3$}
\bkpans{（1）\nobreak $x=\pm5j$。（2）\nobreak $x=\pm\sqrt3\,j$。}

\bkproblem 次の2次方程式の判別式 $D$ を求め、解の種類を答えてから、解きましょう\mbox{（\bkreidairef{8}）}。\newline\bkviiit[0.4]{（1）$s^2+7s+12=0$}\bkviiit[0.4]{（2）$s^2+2s+10=0$}
\bkpans{（1）\nobreak $D=49-48=1>0$。実数の解2つ。$s=-3$、$-4$。（2）\nobreak $D=4-40=-36<0$。虚数の解2つ。$s=-1\pm3j$。}

\bkproblemgroup{チャレンジ}{4}

\bkproblem $s^2+6s+10=0$ を解き、2つの解の和と積を求めて、式の係数と合うか確かめましょう\mbox{（\bkreidairef{7}）}。
\bkpans{$D=36-40=-4$。$s=\dfrac{-6\pm2j}{2}=-3\pm j$。和 $-6$、積 $9+1=10$ で、合っている。}

\bkproblem $\sqrt{-3}\times\sqrt{-12}$ を計算しましょう。また、$\sqrt{(-3)\times(-12)}$ とくらべましょう\mbox{（\bkreidairef{6}）}。
\bkpans{$\sqrt3\,j\times\sqrt{12}\,j=\sqrt{36}\,j^2=-6$。$\sqrt{(-3)\times(-12)}=\sqrt{36}=6$ で、ちがう値になる。}

\bkproblem たすと4、かけると5になる2つの数を求めましょう。また、たすと4、かけると5になることを確かめましょう\mbox{（\bkreidairef{9}）}。
\bkpans{$x(4-x)=5$ より $x^2-4x+5=0$。$D=16-20=-4$、$x=\dfrac{4\pm2j}{2}=2\pm j$。2つの数は $2+j$ と $2-j$。和4、積 $4+1=5$。}

\begin{furikaeri}
\begin{checklist}
\item $j$ をかけると90°回ることから、$j^2=-1$ を説明できる。\bkfurref{bk:6-8:kaiten}
\item 複素数のたし算・ひき算・かけ算・わり算ができる。\bkfurref{bk:6-8:kakeru}
\item 判別式がマイナスの2次方程式を解いて、虚数の解を求められる。\bkfurref{bk:6-8:kai}
\end{checklist}
\end{furikaeri}

\bkfukidashi{sora}{「ない」と思っていた数は、平面の上にあったんだね。次の章では、複素数のかけ算が「回転と拡大」になるしくみを、平面の上でくわしく見よう！}
